\lecture{1}
\title{Introduction, Finite Automata, Regular Expressions}

\begin{document}

The class website is
\href{https://math.mit.edu/~sipser/18404/}{https://math.mit.edu/~sipser/18404/},
taught by Michael Sipser.
This course follows the textbook
\emph{Introduction to the Theory of Computation, 3rd edition}
(also written by Michael Sipser).

\subsection{Finite Automata}

One simplified model of computation is the \textbf{finite automaton}.
Here's a prototypical example.
\begin{center}
    \frame{\includegraphics[width=0.5\textwidth]{lecture-01/automata-example}}
\end{center}
\textbf{Definition.} A finite automaton $M = (Q, \Sigma, \delta, q_0, F)$
is a five-tuple of the following:
\begin{itemize}
    \item A finite set $Q$ of \textbf{states}.  In this example,
    we have $Q = \{q_1, q_2, q_3\}$.
    \item A finite \textbf{input alphabet} $\Sigma$.  In this example,
    we have $\Sigma = \{\str{0}, \str{1}\}$.
    \item A \textbf{transition function} $\delta \colon Q \times \Sigma \to Q$.
    In this example, the map $\delta$ is given
    by the labels of the edges of a directed graph.
    \item A \textbf{start state} $q_0 \in Q$.  In this example,
    we have $q_0 = q_1$.
    \item A set $F \subseteq Q$ of \textbf{accept states}.  In this example,
    we have $F = \{q_3\}$.
\end{itemize}
We say that ``$M$ \textbf{accepts} $w = w_1w_2\dots w_n$''
if starting at $q_0$ and following the arrows according to $w$
lands at an accept state.

\textbf{Example.}
In the example above,
we say that $M$ accepts \str{01101} but rejects \str{00101}.

\begin{remark}
    Note that an accept state of a finite automaton
    need not be ``stable'';
    it may have arrows to other states.
\end{remark}

\subsection{(Regular) Languages}

\textbf{Definition.}
Given a finite automaton $M$,
the \textbf{language} $L(M)$ of $M$
is the set of all (finite) strings that $M$ accepts.
Furthermore, we say that
``$M$ \textbf{recognizes} $A$'' for some language $A$
if and only if $A = L(M)$.

\textbf{Example.}
In the example above,
the language $L(M)$ of $M$ is
the set of (finite) strings with \str{11} as a substring.

\textbf{Definition.}
A language $A$ is \textbf{regular}
if $A$ is the language of some finite automaton.

\textbf{Example.} It turns out that\dots
\begin{itemize}
    \item The language of all strings with an even number of \str{1}s
    is regular.  (Obviously.)
    \item The language of all strings with an equal number of
    \str{0}s and \str{1}s is not regular.  (Not so obviously.)
\end{itemize}
Our goal is to be able to identify what the regular languages are.

\subsection{Regular Operations and Regular Expressions}

It helps to define the following:

\textbf{Definition.}
For any two languages $A$ and $B$,
we can construct new languages via the \textbf{regular operations}:
\begin{itemize}
    \item Their \textbf{union} is
    $A \cup B := \{w \mid w \in A \text{ or } w \in B\}$.
    \item Their \textbf{concatenation} is
    $A \circ B = AB = \{xy \mid x \in A, y \in B\}$.
    \item The \textbf{star} of a language is
    $A^* := \{x_1x_2\dots x_k \mid k \geq 0, x_i \in A\}$
    (that is, the regex \verb|*|).
    \begin{itemize}
        \item Note that $A^*$ always contains
        the \emph{empty string}, which we denote $\epsilon$.
    \end{itemize}
\end{itemize}

\textbf{Example.}
If $A = \{\str{01}, \str{10}\}$ and
$B = \{\str{22}, \str{33}, \str{44}\}$,
then we have:
\[ A \cup B = \{\str{01}, \str{10},
\str{22}, \str{33}, \str{44}\}
~~~ \text{ and }  ~~~
A \circ B = \{\str{0122}, \str{0133}, \str{0144},
\str{1022}, \str{1033}, \str{1044}\}. \]
Meanwhile, we have $A^* = \{\epsilon,
\str{01}, \str{10}, \str{0101}, \str{0110}, \str{1001}, \str{1010},
\str{010101}, \str{010110}, \dots\}$.

\textbf{Definition.}
We say that a \textbf{regular expression}
is any expression formed using $\epsilon$, $\emptyset$,
atoms $a \in \Sigma$, and regular operations.
\begin{itemize}
    \item Write $(\str{0} \cup \str{1})^*$
    to mean the set of all finite (possibly empty) binary strings.
    \item Write $(\str{0} \cup \str{1})^*\str{1}$
    to mean the set of all binary strings that end in \str{1}.
    \item Write $(\str{0} \cup \str{1})^*\str{11}(\str{0} \cup \str{1})^*$
    to mean the set of all binary strings with \str{11} as a substring.
    \item Things like
    $(\str{0} \cup \str{1})^*\str{11}(\str{0} \cup \str{1})^*
    \cup (\str{0} \cup \str{1})^*\str{00}(\str{0} \cup \str{1})^*$
    are also regular expressions.
\end{itemize}
As the name suggests, the following turns out to be true.

\begin{theorem}{Regular Expressions are Regular}
    A language is regular if and only if some regular expression describes it.
\end{theorem}

For the rest of lecture,
we'll work through a proof of one step of the ``if'' direction.

\begin{theorem}{Union of Regular Languages}
    If $A_1$ and $A_2$ are regular languages,
    then so is $A_1 \cup A_2$.
\end{theorem}

\begin{proof}
    The idea is to use the automata $M_1$ and $M_2$ that recognize
    $A_1$ and $A_2$ (respectively)
    and run them both simultaneously in a single automaton
    by taking their ``cross product''.
    Here's an example of what that might look like.
    \begin{center}
        \frame{\includegraphics[width=0.6\textwidth]{lecture-01/union-proof}}
    \end{center}
    Here, $M_1$ and $M_2$ are depicted as
    ``one-dimensional'' automata,
    drawn horizontally and vertically, respectively.
    Their cross product $M$ is the $3 \times 3$ grid of states
    above and to the right of them.
    Evidently, the language of $M$ is $A_1 \cup A_2$.

    Formally, writing $M_k = (Q_k, \Sigma_k, \delta_k, q_{k, 0}, F_k)$,
    the proof is to consider
    the finite automaton $M = (Q, \Sigma, \delta, q_0, F)$ given by:
    \begin{itemize}
        \item The state set $Q := Q_1 \times Q_2
        = \{(q_{1, i}, q_{2, j}) \mid
        q_{1, i} \in Q_1, q_{2, j} \in Q_2\}$.
        \item The input alphabet $\Sigma := \Sigma_1 = \Sigma_2$.
        \item The transition function $\delta((q_{1, i}, q_{2, j}), a)
        := (\delta_1(q_{1, i}, a), \delta_2(q_{2, j}, a))$.
        \item The start state $q_0 := (q_{1, 0}, q_{2, 0})$.
        \item The accept states $F := \{(q_{1, i}, q_{2, j})
        \mid q_{1, i} \in F_1 \text{ or } q_{2, j} \in F_2\}$.
    \end{itemize}
    And that's it. ~ $\blacksquare$
\end{proof}

\begin{remark}
    Implicitly, we've assumed that the languages $A_1$ and $A_2$
    use the same alphabet $\Sigma = \Sigma_1 = \Sigma_2$.
    We will almost always assume this in general,
    but even without this assumption,
    the argument is easily repairable by
    taking $\Sigma := \Sigma_1 \cup \Sigma_2$ instead.
\end{remark}

By replacing ``$F := \{(q_{1, i}, q_{2, j})
\mid q_{1, i} \in F_1 \text{ or } q_{2, j} \in F_2\}$''
with ``$F := F_1 \times F_2$'' in the proof,
one can also prove the following:

\begin{theorem}{Intersection of Regular Languages}
    If $A_1$ and $A_2$ are regular languages,
    then so is $A_1 \cap A_2$.
\end{theorem}

\end{document}